\documentclass[english,twoside, openright, hidelinks, 11 pt, class=report,crop=false]{standalone}
\input{../../preamb}
\input{../bm}

\begin{document}	
\subsection*{Kapittel \ref{Vek}}

\opr{opgvekfinvek}
$\vv{AB}=[1, 3]$, $\vv{AC}=[-4, -5]$, $\vv{BC}=[-5, -8]$ og $\vv{CD}=[9, -2]$


\opr{opgvekfindvek2}
\selos

\opr{opgvekmidt}
\selos

\opr{opgvekskalar}
$\vec{u}\cdot\vec{v}=32$ og $\vec{u}\cdot\vec{w}=4$

\opr{opgvekskalargeneral}
\selos

\opr{opgveklen}
$ |\vec{a}| = 5
|\vec{b}| = \sqrt{50},
|\vec{c}| = 10,
|\vec{d}| = 5 $

\opr{vekfaktlen}
\selos


\opr{opgvekvislength}
\selos



\opr{opgveknorm}
$\vec{a}$ og $\vec{c}$ er ortogonale. $\vec{a}$ og $\vec{d}$ er ortogonale.

\opr{opgvekskalar2}
Finn skalarproduktet av $\vec{u}$ og $\vec{v}$ når 
\abch{
	\item $\vec{u}\cdot\vec{v}=10\sqrt{3}$
	\item $\vec{u}\cdot\vec{v}=\frac{21\sqrt{2}}{2}$
	\item $\vec{u}\cdot\vec{v}=-\frac{5}{2}$
}

\opr{opgvekrett}
$[a, b]\cdot[-b, a]=-ab+ab=0$

\opr{opgvekcalc}
\abch{
	\item $-4$
	\item $41$
}

\opr{opgvekfindangle}
Finn verdien til $\angle(\vec{u}, \vec{v})$ når du vet at
\abc{
	\item $30^\circ$
	\item $60^\circ$
	\item $180^\circ$
}

\opr{opgveka} $\vec{u}\cdot\vec{v}\leq|\vec{u}||\vec{v}|$

\opr{opgvekparl} Ingen er parallelle.


\opr{opgvekin}
\selos

\opr{opgvekfindsandr}
\abch{
	\item $(3, -1)$ og $4$
	\item $(0, 5)$ og $9$	
}

\opr{opgvekvekfunk}
\abch{
	\item $(2, 3)$ og $[4, -5]$
	\item $-6$
	\item Begge ligger på linja.
}

\opr{opgvekfindfunk}
$\vec{l}(t)=[-2, 7]+t[5, -12]$, $\vec{l}(t): \left\lbrace{
	\begin{array}{l}
		x= -2+5t\\
		y= 7-12t   
	\end{array}
}\right. $ 


\opr{opgvekstigning}
$\frac{b}{a}$

\opr{vekskalparl}
\selos

\grubr{R1V23D1opg3} 
\abch{
	\item  $\angle CBA=90^\circ$
	\item $P=(4, 6)$
}

\grubr{opgvektenh}
\abch{
	\item \selos
	\item $[6, 8]$
}

\grubr{opgveksirkline}
\abch{
	\item $(0, 4)$ og $5$
	$\vec{a}$ skjærer i to punkt, $\vec{b}$ tangerer, $\vec{c}$ skjærer ikke.
	\item Tangeringspunkt: $(-4, 7)$, skjæringspunkt: $(0,-1)$ og $(4, 7)$
}


\grubr{r1h23d1opg3}

\abch{
	\item $BC$ 
	\item Ingen
}

\grubr{r1v22d2opg2}
\abch{
	\item $|\vec{u}|=\sqrt{7}$, $\vec{v}=\sqrt{364}$
	\item $-\frac{35}{\sqrt{7}\sqrt{364}}$
}

\grubr{r1v24d1opg4}
\abch{
	\item $t=2$
	\item $t=0\vee t=-4$	
}

\grubr{r1v22d1opg4}
\abch{
	\item $t=6$
	\item $t=3$
}

\grubr{r1h25d1opg4}
\abch{
	\item $\sqrt{26}$
	\item $(4,0 )$
	\item $t=5$
} 

\grubr{opgvekdiag}
\selos


\grubr{opgveklinperp}
\selos

\grubr{r1h22d1opg4}
\abch{
	\item \selos
	\item $Q=\left(\frac{1}{5}, \frac{27}{5}\right)$
}


\grubr{r1h22d2opg3}
\abch{
	\item \selos
	\item $\frac{226}{26\sqrt{82}}$
	\item \selos
	\item $E=\left(8, \frac{10}{3}\right)$
}

\grubr{opgveksirkeltangent}
$[-\sqrt{3}, 1]$


\grubr{r1h24d1opg5} $\vec{p}=[8, 12]$
\abch{
	\item Nei.
	\item $\vec{u}$ og $\vec{v}$ er ortogonale, $\vec{u}$ og $\vec{p}$ er ortogonale
	\item $a=2$, $b=1$
}

\grubr{r1h25d2opg5}
\abch{
	\item $|\vec{p}|=2\sqrt{10}$
	\item $t=-1$
}

\grubr{r1v25d1opg6} (R1V25) 
\abch{
	\item $AB=\sqrt{5}$
	\item $a=1$
	\item $(1, 3)$
}

\grubr{r1h25d1opg4}
\abch{
	\item $\sqrt{26}$
	\item $(13, 0)$
	\item $t=-3$
}

\grubr{r1v23d2opg2}
\abch{
	\item $C=(7+2t, 5+5t)$, $D=(3+2t, 2+5t)$
	\item $t=3$
}
\end{document}